\chapter{Wittig en andere reacties die C=C vormen}\label{ch:b2:wittig}

Een vrouwelijke huisvlieg lokt mannetjes met een koolwaterstof van 23 koolstofatomen met één enkele dubbele binding, tussen
koolstofatoom 9 en 10, met de (\textit{Z})-geometrie. Om die in grote hoeveelheid te maken, moet een chemicus die ene
dubbele binding precies daar zetten, en met die geometrie: een eliminatie zou haar over verscheidene
posities verspreiden en beide geometrieën geven. Dit hoofdstuk bouwt dubbele bindingen op de andere manier, door twee
stukken te verbinden aan een carbonylgroep, zodat de dubbele binding verschijnt waar de \ce{C=O} zat.
% ledger: use:muscalure

\begin{recall}
\Cref{ch:b2:enolates-aldol}: de \omterm{def:b2:enolates-aldol:aldol}{aldolcondensatie}. \Cref{ch:b2:alkene-redox}: de hydrogenering van
een alkyn tot een (\textit{Z})-alkeen over een vergiftigde katalysator. Het deel van jaar~1: E1, E2 en de regel van
Zaitsev, bimoleculaire substitutie, grignardreagentia en polariteitsinversie, de
descriptoren (\textit{Z})/(\textit{E}). Het schoolboek: atoomeconomie.
\end{recall}

\section{Fosfoniumzouten en ylides}

\begin{definition}[Fosfoniumylides]\label{def:b2:wittig:ylide}
Een \emph{fosfoniumzout}\index{fosfoniumzout} $\mathrm{R_3P^+{-}CH_2R'\ X^-}$ wordt gemaakt door bimoleculaire
substitutie van een halogeenalkaan met een fosfine, meestal trifenylfosfine \ce{Ph3P}. Wie een
waterstofatoom weghaalt van het koolstofatoom dat aan fosfor gebonden is, krijgt een \emph{fosfoniumylide}\index{fosfoniumylide},
een neutraal molecuul met tegengestelde ladingen naast elkaar, $\mathrm{Ph_3P^+{-}CH^-{-}R'}$, ook geschreven als $\mathrm{Ph_3P{=}CHR'}$. Een
\emph{gestabiliseerd ylide}\index{gestabiliseerd ylide} draagt op dat koolstofatoom een elektronenzuigende groep
(een ester, een keton, een \omterm{def:b2:amines:nitrile}{nitril}) die de negatieve lading delokaliseert.
\end{definition}

\begin{omfigure}
\setchemfig{atom sep=1.9em}
\schemestart
\chemfig{Ph_3P} \+ \chemfig{CH_3-I}
\arrow{->[S$_\mathrm{N}$2]}[,0.9]
\chemfig{Ph_3P^{+}-CH_3}\;\chemfig{I^{-}}
\arrow{->[\ce{BuLi}][$-$ \ce{BuH}, \ce{LiI}]}[,1.3]
\chemfig{Ph_3P^{+}-\chemabove{C}{\scriptstyle -}H_2}
\arrow{<->}[,0.8]
\chemfig{Ph_3P=CH_2}
\schemestop
\par\vspace{6pt}

{\small Methylideentrifenylfosforaan, het eenvoudigste ylide. Trifenylfosfine verdringt jodide
uit joodmethaan; butyllithium haalt een proton weg van de methylgroep. De twee structuren rechts
zijn twee schrijfwijzen van hetzelfde molecuul: die met gescheiden ladingen ligt het dichtst bij de werkelijkheid.}
\end{omfigure}

\begin{proposition}[Het ylide maken]\label{prop:b2:wittig:ylide-acidity}
Een \omterm{def:b2:wittig:ylide}{fosfoniumzout} is op het koolstofatoom naast fosfor veel zuurder dan een alkaan. Een eenvoudig
alkylfosfoniumzout vraagt nog steeds een zeer sterke base (butyllithium, natriumhydride, natriumamide) onder droge,
inerte omstandigheden; een zout waarvan de \ce{CH2} ook een carbonylgroep draagt, wordt gedeprotoneerd door een zwakke base
(natriumcarbonaat, natriumhydroxide), en het gevormde \omterm{def:b2:wittig:ylide}{gestabiliseerde ylide} kan geïsoleerd en bewaard worden.
\end{proposition}

\begin{proof}[Redenering]
Het positieve fosforatoom naast het carbanion stabiliseert het door zijn lading en door de polariseerbaarheid van
het grote fosforatoom: een eenvoudig ylide is een gestabiliseerd carbanion, maar minder dan een \omterm{def:b2:enolates-aldol:enolate}{enolaat},
vandaar de sterke base. Een naburige carbonylgroep spreidt de negatieve lading naar zuurstof, zoals in een \omterm{def:b2:enolates-aldol:enolate}{enolaat}
(\cref{prop:b2:enolates-aldol:alpha-acidity}), bovenop het fosforatoom: het geconjugeerde zuur wordt
zuur genoeg voor een zwakke base, en het ylide, veel minder basisch, reageert niet meer met water of lucht.
\end{proof}

\section{De wittigreactie}

\begin{definition}[Wittigreactie]\label{def:b2:wittig:wittig}
De \emph{wittigreactie}\index{wittigreactie} van een \omterm{def:b2:wittig:ylide}{fosfoniumylide} met een aldehyde of een keton
geeft een alkeen en trifenylfosfineoxide:
\begin{equation*}
\mathrm{Ph_3P{=}CHR + R'CHO \longrightarrow R'CH{=}CHR + Ph_3P{=}O}.
\end{equation*}
Ze verloopt via een \emph{oxafosfetaan}\index{oxafosfetaan}, een vierring van fosfor, twee
koolstofatomen en zuurstof, die ontstaat doordat het ylidekoolstofatoom aan het carbonylkoolstofatoom bindt terwijl het carbonylzuurstofatoom
aan fosfor bindt.
\end{definition}

\begin{omfigure}
\setchemfig{atom sep=1.9em}
\schemestart
\chemfig{Ph_3P=CH-R} \+ \chemfig{O=CH-R'}
\arrow{->[additie]}[,1.0]
\chemfig{Ph_3P?-[6]CH(-[4]R)-[0]CH(-[0]R')-[2]O?}
\schemestop

\medskip
\schemestart
\arrow{->[eliminatie]}[,1.1]
\chemfig{R-CH=CH-R'} \+ \chemfig{Ph_3P=O}
\schemestop
\par\vspace{6pt}

{\small De \omterm{def:b2:wittig:wittig}{wittigreactie}. Het ylidekoolstofatoom en het carbonylkoolstofatoom binden aan elkaar terwijl zuurstof aan
fosfor bindt, in één stap: het \omterm{def:b2:wittig:wittig}{oxafosfetaan}. De ring gaat daarna de andere kant op open, waarbij de bindingen
\ce{P-C} en \ce{C-O} breken: de binding \ce{C=C} vormt zich tussen de twee vroegere reagerende koolstofatomen, en de
binding \ce{P=O} van trifenylfosfineoxide vormt zich tegelijk.}
\end{omfigure}

\begin{proposition}[Een dubbele binding precies op haar plaats]\label{prop:b2:wittig:regiospecific}
Bij een \omterm{def:b2:wittig:wittig}{wittigreactie} verbindt de nieuwe binding \ce{C=C} het vroegere carbonylkoolstofatoom met het vroegere ylidekoolstofatoom,
en nergens anders.
\end{proposition}

\begin{proof}[Redenering]
De enige bindingen die gevormd en verbroken worden, zijn die van de vierring: het ylidekoolstofatoom aan het
carbonylkoolstofatoom, daarna de bindingen \ce{C-P} en \ce{C-O}. Geen enkel waterstofatoom verschuift en er vormt zich geen carbokation of carbanion
langs de keten, dus de dubbele binding kan niet verschuiven. Een eliminatie kiest daarentegen tussen de
waterstofatomen op verscheidene naburige koolstofatomen (regel van Zaitsev), en een E1-reactie kan omleggen.
\end{proof}

\begin{proposition}[Drijvende kracht]\label{prop:b2:wittig:driving}
De vorming van de binding \ce{P=O} van trifenylfosfineoxide maakt de \omterm{def:b2:wittig:wittig}{wittigreactie} sterk
gunstig en onomkeerbaar.
\end{proposition}

\begin{proof}[Redenering]
De reactie ruilt een binding \ce{C=O} en een binding \ce{P-C} (in het ylide is de \ce{P=C} grotendeels een
enkelvoudige binding \ce{P+-C-}) in voor een binding \ce{C=C} en een binding \ce{P=O}. De binding \ce{P=O} van een fosfineoxide is
zeer sterk, veel sterker dan de verloren binding \ce{P-C}, en de verloren \ce{C=O} wordt gedeeltelijk gecompenseerd
door de gevormde \ce{C=C}: de balans gaat duidelijk bergaf. Fosfor is de zuurstofput van de
reactie.
\end{proof}

\section{Stereoselectiviteit}

\begin{proposition}[Geometrie van het alkeen]\label{prop:b2:wittig:stereo}
Met een aldehyde geeft een niet-gestabiliseerd ylide (\ce{Ph3P=CHR}, R een alkylgroep) vooral het (\textit{Z})-alkeen,
onder zoutvrije omstandigheden. Een \omterm{def:b2:wittig:ylide}{gestabiliseerd ylide} (\ce{Ph3P=CH-COOR}, \ce{Ph3P=CH-COR}) geeft vooral het
(\textit{E})-alkeen. Een ylide dat alleen door een arylgroep gestabiliseerd wordt, geeft mengsels.
\end{proposition}

\admitted

Dit zijn waarnemingen. De gebruikelijke verklaring leest de \omterm{def:b2:wittig:wittig}{wittigreactie} als een competitie tussen
snelheden en stabiliteiten. Met een reactief, niet-gestabiliseerd ylide vormt het \omterm{def:b2:wittig:wittig}{oxafosfetaan} zich snel en
onomkeerbaar, via de minst gedrongen overgangstoestand, waarin de twee substituenten aan dezelfde
kant van de ring terechtkomen; het valt daarna uiteen met behoud van die schikking: kinetische controle geeft het
(\textit{Z})-alkeen. Met een \omterm{def:b2:wittig:ylide}{gestabiliseerd ylide} is de eerste stap trager en omkeerbaar; de
oxafosfetanen gaan via de uitgangsstoffen in elkaar over, en dat met de substituenten aan
tegenovergestelde kanten, minder gespannen, wint: het (\textit{E})-alkeen, onder thermodynamische controle.

\begin{omfigure}
\setchemfig{atom sep=1.9em}
\schemestart
\chemfig{Ph_3P=CH-CH_2CH_3} \+ \chemfig{O=CH-Ph}
\arrow{->}[,0.8]
\chemfig{Ph-[:-30]=[:30]-[:90]CH_2CH_3}
\schemestop

\medskip
\schemestart
\chemfig{Ph_3P=CH-COOEt} \+ \chemfig{O=CH-Ph}
\arrow{->}[,0.8]
\chemfig{Ph-[:-30]=[:30]-[:-30]COOEt}
\schemestop
\par\vspace{6pt}

{\small Twee ylides, één aldehyde. Boven: het niet-gestabiliseerde propylideenylide geeft vooral
(\textit{Z})-1-fenylbut-1-een. Onder: het \omterm{def:b2:wittig:ylide}{gestabiliseerde ylide} geeft vooral
ethyl-(\textit{E})-3-fenylpropenoaat (ethylcinnamaat). Trifenylfosfineoxide, dat in beide gevormd wordt, is niet
getekend.}
\end{omfigure}

\section{De Horner--Wadsworth--Emmons-reactie}

\begin{definition}[Fosfonaten en de HWE-reactie]\label{def:b2:wittig:hwe}
Een \emph{fosfonaat}\index{fosfonaat} is een ester van een fosfonzuur, $\mathrm{(RO)_2P({=}O){-}R'}$, met een
binding \ce{P-C}. Het wordt gemaakt door de Michaelis--Arbuzov-reactie van een trialkylfosfiet met een
halogeenalkaan, $\mathrm{P(OEt)_3 + R'Br \to (EtO)_2P(O)R' + EtBr}$. Bij de
\emph{Horner--Wadsworth--Emmons-reactie}\index{Horner--Wadsworth--Emmons-reactie} (HWE) addeert het anion van een
fosfonaat met een elektronenzuigende groep op zijn \ce{CH2} aan een aldehyde, en geeft het een
alkeen en een dialkylfosfaatanion.
\end{definition}

\begin{omfigure}
\setchemfig{atom sep=1.8em}
\schemestart
\chemfig{EtO-P(=[2]O)(-[6]OEt)-CH_2-COOEt}
\arrow{->[\ce{NaH}][$-$ \ce{H2}]}[,0.9]
\chemfig{EtO-P(=[2]O)(-[6]OEt)-\chemabove{C}{\scriptstyle -}H-COOEt}
\schemestop

\medskip
\schemestart
\arrow{->[\ce{PhCHO}]}[,0.9]
\chemfig{*6(=-=(-[:30]=[:-30]-[:30](=[:90]O)-[:-30]O-[:30]Et)-=-)}
\+
\chemfig{EtO-P(=[2]O)(-[6]O^{-})-OEt}
\schemestop
\par\vspace{6pt}

{\small De \omterm{def:b2:wittig:hwe}{Horner--Wadsworth--Emmons-reactie} van triethylfosfonoacetaat met benzaldehyde.
Natriumhydride haalt het waterstofatoom tussen het \omterm{def:b2:wittig:hwe}{fosfonaat} en de ester weg; het anion addeert aan het aldehyde,
en het vierledige intermediair valt uiteen zoals bij de \omterm{def:b2:wittig:wittig}{wittigreactie}. Het product is ethyl-(\textit{E})-cinnamaat;
het bijproduct is het diethylfosfaatanion, oplosbaar in water.}
\end{omfigure}

\begin{proposition}[Waarom chemici de HWE-reactie waarderen]\label{prop:b2:wittig:hwe}
De HWE-reactie geeft vooral het (\textit{E})-alkeen, en haar fosforbijproduct, een in water oplosbaar
zout, wordt verwijderd door eenvoudig met water te wassen; het fosfonaatanion is bovendien nucleofieler dan het
overeenkomstige \omterm{def:b2:wittig:ylide}{gestabiliseerde ylide} en reageert ook met ketonen.
\end{proposition}

\begin{proof}[Redenering]
Het fosfonaatanion is een gestabiliseerd carbanion, dus zijn additie is omkeerbaar en de reactie verloopt
onder thermodynamische controle: (\textit{E}), zoals bij \omterm{def:b2:wittig:ylide}{gestabiliseerde ylides}. Anders dan een \omterm{def:b2:wittig:ylide}{fosfoniumylide}
draagt het een volledige negatieve lading, die niet door een naburig kation gecompenseerd wordt: het is reactiever. Het
bijproduct, \ce{(EtO)2PO2-}, is een ionisch zout, terwijl trifenylfosfineoxide een neutrale organische
vaste stof is die door chromatografie of kristallisatie van het product gescheiden moet worden.
\end{proof}

\section{Een methode kiezen}

\begin{method}[Wittigdisconnectie]\label{met:b2:wittig:wittig-disconnection}
\begin{enumerate}
  \item Knip de beoogde binding \ce{C=C} door; het ene koolstofatoom wordt een carbonylkoolstofatoom, het andere het ylidekoolstofatoom
        (een \omterm{def:b2:wittig:ylide}{fosfoniumzout}, uit een halogenide).
  \item Verkies van de twee mogelijkheden die waarvan het halogenide primair is (bimoleculaire substitutie door
        \ce{Ph3P}) en waarvan de carbonylverbinding een aldehyde is.
  \item Kies het reagens volgens de gewenste geometrie: een niet-gestabiliseerd ylide voor (\textit{Z}), een
        \omterm{def:b2:wittig:ylide}{gestabiliseerd ylide} of een \omterm{def:b2:wittig:hwe}{fosfonaat} (HWE) voor (\textit{E}).
  \item Controleer dat niets anders in het molecuul met de gebruikte base reageert.
\end{enumerate}
\end{method}

\begin{method}[Een reactie kiezen die C=C vormt]\label{met:b2:wittig:choose-cc}
\begin{center}
\small
\begin{tabular}{@{}lll@{}}
\hline
Methode & Positie van de \ce{C=C} & Geometrie \\
\hline
E2- of E1-eliminatie & regel van Zaitsev, mengsels & meestal (\textit{E}), mengsels \\
\omterm{def:b2:enolates-aldol:aldol}{Aldolcondensatie} & geconjugeerd met \ce{C=O} & (\textit{E}) \\
Wittig, niet-gestabiliseerd ylide & op de vroegere \ce{C=O} & (\textit{Z}) \\
Wittig, \omterm{def:b2:wittig:ylide}{gestabiliseerd ylide}; HWE & op de vroegere \ce{C=O} & (\textit{E}) \\
Alkyn, dan vergiftigde katalysator & op de vroegere drievoudige binding & (\textit{Z}) \\
\hline
\end{tabular}
\end{center}
\medskip
Om twee stukken van vergelijkbare grootte te verbinden aan een \ce{C=C} zijn de \omterm{def:b2:wittig:wittig}{wittigreactie} en haar verwanten meestal
de eerste keuze; olefinemetathese, die de uiteinden van twee alkenen uitwisselt, komt in het deel
van jaar~3.
\end{method}

\begin{history}[Het ylide dat een reagens werd]
Georg Wittig ontdekte begin jaren 1950, toen hij fosforylides bestudeerde, dat methylideentrifenylfosforaan
benzofenon omzet in 1,1-difenyletheen en trifenylfosfineoxide; met zijn student Ulrich
Schöllkopf maakte hij van die waarneming in 1954 een algemene methode. De industrie nam ze snel over, onder meer voor
vitamine A, en Wittig deelde in 1979 de Nobelprijs voor de Scheikunde.
\end{history}

\begin{safety}{\ghs[5mm]{GHS02}\,\ghs[5mm]{GHS05}\,\ghs[5mm]{GHS07}\,\ghs[5mm]{GHS08}}
Butyllithiumoplossingen zijn ontvlambaar en corrosief en reageren heftig met water; natriumhydride
maakt met water waterstof vrij. Trifenylfosfine is schadelijk en corrosief voor de ogen,
triethylfosfiet ontvlambaar en schadelijk; trifenylfosfineoxide is schadelijk.
% ledger: ghs:BuLi, ghs:PPh3, ghs:triethyl-phosphite, ghs:Ph3PO
\end{safety}

\section{Oefeningen}

\begin{exercise}[$\star$]\label{exo:b2:wittig:1}
Schrijf de bereiding op van het ylide \ce{Ph3P=CHCH3} uit trifenylfosfine en een halogeenalkaan.
\end{exercise}

\begin{exercise}[$\star$]\label{exo:b2:wittig:2}
Geef de producten van de \omterm{def:b2:wittig:wittig}{wittigreactie} van benzaldehyde met het ylide dat uit joodmethaan gemaakt wordt.
\end{exercise}

\begin{exercise}[$\star$]\label{exo:b2:wittig:3}
Deel in als gestabiliseerd, niet-gestabiliseerd of arylgestabiliseerd: \ce{Ph3P=CHCOOEt}, \ce{Ph3P=CHCH3},
\ce{Ph3P=CHCOCH3}, \ce{Ph3P=CHC6H5}. Welk ervan vraagt butyllithium?
\end{exercise}

\begin{exercise}[$\star$]\label{exo:b2:wittig:4}
Geef het product van triethylfosfonoacetaat, natriumhydride en propanal, met zijn geometrie.
\end{exercise}

\begin{exercise}[$\star\star$]\label{exo:b2:wittig:5}
Men wil methyleencyclohexaan. Vergelijk een wittigroute uit cyclohexanon met de zuurgekatalyseerde
dehydratatie van 1-methylcyclohexaan-1-ol.
\end{exercise}

\begin{exercise}[$\star\star$]\label{exo:b2:wittig:6}
Maak een wittigdisconnectie van (\textit{E})-1,2-difenyletheen (stilbeen). Waarom is de wittigroute
een slechte keuze voor zuiver (\textit{E})-stilbeen, en wat zou je in plaats daarvan gebruiken?
\end{exercise}

\begin{exercise}[$\star\star$]\label{exo:b2:wittig:7}
Propanal reageert enerzijds met \ce{Ph3P=CHCH3} en anderzijds met \ce{Ph3P=CHCOOEt}. Geef de
producten en hun voornaamste geometrie.
\end{exercise}

\begin{exercise}[$\star\star$]\label{exo:b2:wittig:8}
Bereken de atoomeconomie van de methylenering van cyclohexanon door \ce{Ph3P=CH2}.
\end{exercise}

\begin{exercise}[$\star\star$]\label{exo:b2:wittig:9}
Schrijf de Michaelis--Arbuzov-reactie van triethylfosfiet met ethylbroomethanoaat op, in twee stappen:
substitutie door fosfor, daarna dealkylering door bromide.
\end{exercise}

\begin{exercise}[$\star\star\star$]\label{exo:b2:wittig:10}
Stel een synthese voor van 1,6-difenylhexa-1,3,5-trieen uit (\textit{E})-but-2-eendial
\ce{OHC-CH=CH-CHO} via twee \omterm{def:b2:wittig:wittig}{wittigreacties}, en zeg welk ylide en hoeveel ervan.
\end{exercise}

\begin{exercise}[$\star\star\star$]\label{exo:b2:wittig:11}
(\textit{E})-4-Fenylbut-3-een-2-on kan gemaakt worden door een \omterm{def:b2:enolates-aldol:aldol}{aldolcondensatie} of door een \omterm{def:b2:wittig:wittig}{wittigreactie}.
Schrijf beide routes op en vergelijk ze.
\end{exercise}

\begin{exercise}[$\star\star\star$]\label{exo:b2:wittig:12}
Stel twee routes voor naar (\textit{Z})-oct-4-een, één via de \omterm{def:b2:wittig:wittig}{wittigreactie} en één via een alkyn,
en zeg welke bijproducten elk ervan achterlaat.
\end{exercise}

\section{Probleem: een feromoon via Wittig}

\begin{problem}[{Weekendprobleem --- de retrosynthese van een (Z)-alkeen, het niet-gestabiliseerde ylide en
zijn selectiviteit, de kost aan trifenylfosfineoxide, en het alternatief via een alkyn}]
\label{pb:b2:wittig:1}
Doel: (\textit{Z})-tricos-9-een, \ce{CH3(CH2)7CH=CH(CH2)12CH3}, het huisvliegferomoon van de
inleiding. Molaire massa's (\unit{g/mol}): C 12.011, H 1.008, O 15.999, P 30.974, Br 79.904.
% ledger: use:muscalure, aw:C, aw:H, aw:O, aw:P, aw:Br, ghs:nonanal, ghs:BuLi

\medskip
\noindent\textbf{Deel I --- Retrosynthese.}
\begin{enumerate}
  \item Schrijf de molecuulformule van het doelmolecuul op en controleer dat het een alkeen is.
  \item Welke binding knipt een wittigdisconnectie door?
  \item Geef de twee mogelijke paren (carbonylverbinding, \omterm{def:b2:wittig:ylide}{fosfoniumzout}).
  \item Welk halogeenalkaan maakt voor het paar nonanal + ylide het \omterm{def:b2:wittig:ylide}{fosfoniumzout}?
  \item Waarom is een primair halogeenalkaan geschikt voor deze stap?
  \item Schrijf de vorming van het \omterm{def:b2:wittig:ylide}{fosfoniumzout} op.
\end{enumerate}

\medskip
\noindent\textbf{Deel II --- Het ylide en de geometrie.}
\begin{enumerate}[resume]
  \item Welke base deprotoneert dit zout? Waarom zou natriumhydroxide niet volstaan?
  \item Is het ylide gestabiliseerd? Welke geometrie voorspelt dat?
  \item Schrijf het \omterm{def:b2:wittig:wittig}{oxafosfetaan} op dat met nonanal gevormd wordt.
  \item Waarom moet de reactie droog en onder inert gas uitgevoerd worden?
  \item Waarom komt de dubbele binding precies tussen koolstofatoom 9 en 10 terecht?
  \item Wat zou de zuurgekatalyseerde dehydratatie van tricosaan-9-ol in plaats daarvan geven?
\end{enumerate}

\medskip
\noindent\textbf{Deel III --- Trifenylfosfineoxide.}
\begin{enumerate}[resume]
  \item Schrijf de totale vergelijking op vanuit het ylide en nonanal.
  \item Bereken de molaire massa's van het doelmolecuul, van nonanal, van het \omterm{def:b2:wittig:ylide}{fosfoniumzout} en van
        trifenylfosfineoxide.
  \item Bereken de atoomeconomie van de wittigstap (ylide + aldehyde).
  \item Hoe wordt trifenylfosfineoxide van een koolwaterstof gescheiden?
  \item Waarom is de vorming van trifenylfosfineoxide de drijvende kracht?
  \item Waarom zou de HWE-reactie niet geschikt zijn voor dit doelmolecuul?
\end{enumerate}

\medskip
\noindent\textbf{Deel IV --- De alkynroute.}
\begin{enumerate}[resume]
  \item Welk alkyn geeft, na hydrogenering, het doelmolecuul, en met welke katalysator?
  \item Bouw dat alkyn op uit dec-1-yn en een halogeenalkaan: welk halogenide, welke base?
  \item Waarom moet de hydrogenering stoppen bij het alkeen, en hoe zorgt de katalysator daarvoor?
  \item Vergelijk de twee routes: aantal stappen, controle over de geometrie, bijproducten.
  \item Wat is de atoomeconomie van de hydrogeneringsstap?
  \item Geef de massa trifenylfosfineoxide die per gram feromoon als nevenproduct ontstaat langs de
        wittigroute.
\end{enumerate}
\end{problem}
